\documentclass[12pt]{article}

\usepackage[utf8]{inputenc}
\usepackage[T1]{fontenc}
\usepackage[romanian]{babel}
\usepackage{amsmath}
\usepackage{amssymb}
\usepackage{array}
\usepackage{geometry}
\usepackage{xcolor}
\usepackage{graphicx}

\geometry{a4paper, margin=2.2cm}
\setlength{\parindent}{0pt}
\setlength{\parskip}{0.65em}

\title{\textbf{Model matematic pentru motorul MiniScript+}}
\author{Andrei Sanke Evergreen}
\date{Mai 2026}

\begin{document}

\geometry{a4paper, margin=2.2cm}
\setlength{\parindent}{0pt}
\setlength{\parskip}{0.65em}

\title{\textbf{Model matematic pentru motorul MiniScript+}}
\author{Andrei Sanke Evergreen}
\date{Mai 2026}

\begin{document}

\maketitle

\section{Introducere}

MiniScript+ este un limbaj educațional proiectat pentru a simplifica procesul de învățare a programării. Din punct de vedere matematic, un program MiniScript+ poate fi privit ca o structură formală alcătuită din instrucțiuni, expresii și o stare internă care se modifică în timpul execuției.

Motorul MiniScript+ funcționează ca un interpretor: primește codul sursă, îl transformă într-o reprezentare internă, apoi execută instrucțiunile pas cu pas. Fluxul principal este:

\[
\text{Cod sursă}
\longrightarrow
\text{Tokenizare}
\longrightarrow
\text{Parsare}
\longrightarrow
\text{AST}
\longrightarrow
\text{Evaluare}
\longrightarrow
\text{Output}
\]

Această structură permite atât rularea programului, cât și vizualizarea execuției, debugging-ul și estimarea complexității algoritmului.

\section{Modelul formal al programului}

Un program MiniScript+ poate fi modelat ca o secvență finită de instrucțiuni:

\[
P = (I_1, I_2, \ldots, I_n)
\]

unde fiecare instrucțiune \(I_i\) aparține unei mulțimi de instrucțiuni admise:

\[
I_i \in \mathcal{I}
\]

Mulțimea instrucțiunilor poate fi descrisă astfel:

\[
\mathcal{I} =
\{
\text{INPUT},
\text{PRINT},
\text{ASSIGN},
\text{IF},
\text{ELSE},
\text{WHILE},
\text{END}
\}
\]

Astfel, un program nu este tratat doar ca text, ci ca o listă ordonată de operații care pot modifica starea internă a interpretorului.

\section{Starea internă a interpretorului}

În timpul execuției, interpretorul păstrează o stare internă. Aceasta poate fi modelată printr-un tuplu:

\[
\sigma = (V, O, l, s, Q)
\]

unde:

\[
\begin{aligned}
V &= \text{tabelul variabilelor definite în program} \\
O &= \text{lista valorilor afișate la output} \\
l &= \text{linia curentă executată} \\
s &= \text{numărul de pași executați} \\
Q &= \text{coada valorilor de input}
\end{aligned}
\]

Tabelul variabilelor poate fi privit ca o funcție parțială:

\[
V : \text{Identificator} \rightarrow \text{Valoare}
\]

De exemplu, dacă în program apare instrucțiunea:

\[
X = 5
\]

atunci starea variabilelor se modifică astfel:

\[
V'(X) = 5
\]

iar pentru orice altă variabilă \(Y \neq X\):

\[
V'(Y) = V(Y)
\]

\section{Funcția de tranziție}

Execuția programului poate fi descrisă printr-o funcție de tranziție:

\[
\delta(\sigma, I_i) = \sigma'
\]

Aceasta înseamnă că, pornind de la o stare curentă \(\sigma\) și executând instrucțiunea \(I_i\), interpretorul obține o stare nouă \(\sigma'\).

Pentru o atribuire de forma:

\[
X = expr
\]

tranziția poate fi descrisă prin:

\[
\delta(\sigma, X = expr) =
(V[X \leftarrow eval(expr, V)], O, l+1, s+1, Q)
\]

unde \(eval(expr, V)\) este valoarea expresiei calculate folosind variabilele curente.

Pentru o instrucțiune de afișare:

\[
PRINT \ expr
\]

tranziția devine:

\[
\delta(\sigma, PRINT \ expr) =
(V, O \cup \{eval(expr, V)\}, l+1, s+1, Q)
\]

Prin urmare, instrucțiunile MiniScript+ pot fi analizate riguros ca transformări asupra unei stări interne.

\section{Tokenizarea}

Tokenizarea este etapa în care codul sursă este parcurs caracter cu caracter și împărțit în unități lexicale numite token-uri.

Dacă notăm codul sursă printr-un șir de caractere:

\[
C = c_1c_2 \ldots c_n
\]

atunci lexerul construiește o secvență de token-uri:

\[
L(C) = (t_1, t_2, \ldots, t_m)
\]

unde fiecare token \(t_i\) poate fi un identificator, un număr, un string, un operator, o paranteză sau un cuvânt cheie.

Exemplu:

\[
\texttt{X = X + 1}
\]

devine:

\[
(\text{IDENT}(X), \text{OP}(=), \text{IDENT}(X), \text{OP}(+), \text{NUMBER}(1))
\]

Deoarece fiecare caracter este citit cel mult de un număr constant de ori, complexitatea tokenizării este:

\[
T_{tokenizare}(n) = O(n)
\]

unde \(n\) este numărul de caractere din codul sursă.

\section{Parsarea și construirea structurii programului}

Parsarea verifică dacă token-urile respectă regulile limbajului și construiește reprezentarea internă a programului.

La nivel simplificat, gramatica limbajului poate fi exprimată astfel:

\[
\begin{aligned}
Program &\rightarrow Instr^* \\
Instr &\rightarrow Assign \mid Print \mid Input \mid If \mid While \\
Assign &\rightarrow id = Expr \\
Print &\rightarrow PRINT \ Expr \\
Input &\rightarrow INPUT \ id \\
If &\rightarrow IF \ Expr \ THEN \ Program \ (ELSE \ Program)? \ END \\
While &\rightarrow WHILE \ Expr \ Program \ END
\end{aligned}
\]

Expresiile sunt analizate ținând cont de prioritatea operatorilor. De exemplu:

\[
X + 2 \cdot Y
\]

trebuie interpretat ca:

\[
X + (2 \cdot Y)
\]

nu ca:

\[
(X + 2) \cdot Y
\]

Parserul citește token-urile și construiește structuri pentru instrucțiuni, condiții, bucle și expresii. Dacă există \(m\) token-uri, iar fiecare token este procesat o singură dată sau de un număr constant de ori, atunci:

\[
T_{parsare}(m) = O(m)
\]

\section{Reprezentarea expresiilor prin AST}

Expresiile MiniScript+ sunt reprezentate printr-un arbore sintactic abstract:

\[
AST_{expr} = (N, E, r)
\]

unde:

\[
\begin{aligned}
N &= \text{mulțimea nodurilor arborelui} \\
E &= \text{mulțimea muchiilor părinte-copil} \\
r &= \text{rădăcina arborelui}
\end{aligned}
\]

Nodurile interne reprezintă operatori, iar frunzele reprezintă valori sau variabile. Pentru expresia:

\[
X + 2 \cdot Y
\]

AST-ul are forma conceptuală:

\[
\begin{array}{c}
+ \\
/ \ \backslash \\
X \quad * \\
\quad / \ \backslash \\
\quad 2 \quad Y
\end{array}
\]

Această reprezentare este importantă deoarece păstrează ordinea corectă de evaluare și prioritatea operatorilor.

\section{Evaluarea recursivă a expresiilor}

Evaluarea unei expresii se face recursiv pe AST. Definim funcția:

\[
eval : AST \times V \rightarrow \text{Valoare}
\]

Pentru o constantă numerică:

\[
eval(c, V) = c
\]

Pentru o variabilă:

\[
eval(x, V) = V(x)
\]

Pentru o expresie binară:

\[
eval(a \ op \ b, V) =
eval(a, V) \ op \ eval(b, V)
\]

De exemplu:

\[
eval(X + 2 \cdot Y, V)
=
V(X) + 2 \cdot V(Y)
\]

Dacă AST-ul unei expresii are \(k\) noduri, fiecare nod este vizitat o singură dată, deci:

\[
T_{evaluare\_expr}(k) = O(k)
\]

Spațiul suplimentar folosit de evaluarea recursivă depinde de înălțimea arborelui:

\[
S_{recursie} = O(h)
\]

unde \(h\) este înălțimea AST-ului.

\section{Execuția instrucțiunilor}

După parsare, programul este executat instrucțiune cu instrucțiune. Dacă notăm prin \(p\) numărul de instrucțiuni efectiv executate, atunci:

\[
T_{executie} = O(p)
\]

Valoarea lui \(p\) nu este întotdeauna egală cu numărul de linii din program, deoarece buclele pot executa aceeași instrucțiune de mai multe ori.

Pentru o buclă:

\[
WHILE \ cond \ DO \ B
\]

dacă blocul \(B\) are costul \(T_B\), iar condiția este adevărată de \(q\) ori, atunci costul total este aproximativ:

\[
T_{while} = O(q \cdot (T_{cond} + T_B))
\]

Pentru două bucle imbricate, cu \(q_1\) și \(q_2\) iterații:

\[
T_{nested} = O(q_1 \cdot q_2 \cdot T_B)
\]

Dacă ambele depind de aceeași dimensiune \(n\), obținem:

\[
T_{nested} = O(n^2)
\]

\section{Limitarea buclelor infinite}

Pentru siguranță, motorul folosește o limită maximă de pași:

\[
s \leq MAX\_STEPS
\]

La fiecare instrucțiune executată, contorul \(s\) crește:

\[
s' = s + 1
\]

Dacă:

\[
s > MAX\_STEPS
\]

execuția este oprită și programul este considerat potențial blocat într-o buclă infinită. Această măsură nu demonstrează matematic existența unei bucle infinite, dar previne blocarea aplicației în browser.

\section{Complexitatea totală a motorului}

Procesarea completă a unui program MiniScript+ include tokenizare, parsare, construire AST și execuție:

\[
T_{total} =
T_{tokenizare} +
T_{parsare} +
T_{AST} +
T_{executie}
\]

Prin înlocuire:

\[
T_{total} = O(n) + O(m) + O(k) + O(p)
\]

deci:

\[
T_{total} = O(n + m + k + p)
\]

unde:

\[
\begin{aligned}
n &= \text{numărul de caractere din codul sursă} \\
m &= \text{numărul de token-uri generate} \\
k &= \text{numărul de noduri din AST} \\
p &= \text{numărul de instrucțiuni executate}
\end{aligned}
\]

Dacă toate aceste mărimi sunt raportate la dimensiunea totală a programului, notată cu \(N\), atunci pentru etapa de procesare statică:

\[
T_{static} = O(N)
\]

Execuția poate deveni mai mare decât \(O(N)\) dacă programul conține bucle care rulează de multe ori.

\section{Complexitatea în spațiu}

Spațiul folosit de motor provine din mai multe structuri interne:

\[
S_{total} =
S_{tokens} +
S_{AST} +
S_{variabile} +
S_{recursie} +
S_{output}
\]

Prin aproximare:

\[
S_{total} = O(m + k + v + h + o)
\]

unde:

\[
\begin{aligned}
m &= \text{numărul de token-uri} \\
k &= \text{numărul de noduri din AST} \\
v &= \text{numărul de variabile distincte} \\
h &= \text{înălțimea maximă a AST-ului} \\
o &= \text{numărul de valori afișate la output}
\end{aligned}
\]

În multe programe educaționale simple, output-ul și numărul de variabile sunt mici, iar spațiul este dominat de token-uri și AST:

\[
S_{total} \approx O(m + k)
\]

\section{Analizatorul de complexitate}

Analizatorul de complexitate nu execută programul pentru toate valorile posibile de input. În schimb, el estimează complexitatea pe baza structurii sintactice a codului. Acesta identifică:

\[
\begin{aligned}
L &= \text{numărul total de bucle} \\
d &= \text{adâncimea maximă a buclelor imbricate} \\
v &= \text{numărul de variabile folosite} \\
b &= \text{numărul de blocuri condiționale}
\end{aligned}
\]

O estimare simplificată pentru complexitatea în timp este:

\[
C_{time} \approx O(n^d)
\]

unde \(d\) este adâncimea maximă a buclelor imbricate detectate.

\[
\begin{array}{|c|c|}
\hline
\textbf{Structura programului} & \textbf{Complexitate estimată} \\
\hline
\text{fără bucle relevante} & O(1) \\
\hline
\text{o buclă simplă} & O(n) \\
\hline
\text{două bucle imbricate} & O(n^2) \\
\hline
\text{trei bucle imbricate} & O(n^3) \\
\hline
\end{array}
\]

Pentru spațiu, analizatorul urmărește variabilele și structurile auxiliare:

\[
C_{space} \approx O(v + a)
\]

unde \(a\) reprezintă eventuale structuri auxiliare detectate. În versiunea educațională a MiniScript+, unde accentul este pus pe variabile simple și control de flux, estimarea spațiului este de obicei:

\[
C_{space} = O(1)
\]

dacă nu există structuri care cresc odată cu input-ul.

\section{Scorul de eficiență}

Pentru a oferi feedback vizual, analizatorul poate transforma informațiile structurale într-un scor de eficiență:

\[
Score = 100 - \alpha L - \beta d - \gamma b
\]

unde:

\[
\begin{aligned}
L &= \text{numărul de bucle} \\
d &= \text{adâncimea maximă a buclelor} \\
b &= \text{numărul de blocuri condiționale} \\
\alpha,\beta,\gamma &= \text{coeficienți de penalizare}
\end{aligned}
\]

Acest scor nu reprezintă o demonstrație formală a performanței, ci o măsură educațională care ajută utilizatorul să observe dacă soluția are multe bucle sau blocuri imbricate.

\section{Exemplu de analiză}

Pentru programul:

\[
\begin{array}{l}
\texttt{X = 0} \\
\texttt{WHILE X < N} \\
\quad \texttt{PRINT X} \\
\quad \texttt{X = X + 1} \\
\texttt{END}
\end{array}
\]

există o singură buclă, deci:

\[
d = 1
\]

Dacă bucla rulează de \(N\) ori, complexitatea în timp este:

\[
T(N) = O(N)
\]

Numărul de variabile folosite este constant, deci:

\[
S(N) = O(1)
\]

Pentru două bucle imbricate:

\[
\begin{array}{l}
\texttt{I = 0} \\
\texttt{WHILE I < N} \\
\quad \texttt{J = 0} \\
\quad \texttt{WHILE J < N} \\
\quad \quad \texttt{PRINT I + J} \\
\quad \quad \texttt{J = J + 1} \\
\quad \texttt{END} \\
\quad \texttt{I = I + 1} \\
\texttt{END}
\end{array}
\]

adâncimea maximă este:

\[
d = 2
\]

iar complexitatea estimată este:

\[
T(N) = O(N^2)
\]

\section{Observații privind corectitudinea analizei}

Analizatorul de complexitate oferă o estimare euristică. El folosește structura codului, nu o demonstrație matematică completă pentru orice input posibil. De exemplu, o buclă de forma:

\[
\texttt{WHILE X < 10}
\]

poate avea un număr constant de iterații, deci complexitatea ei poate fi tratată ca:

\[
O(1)
\]

în timp ce o buclă dependentă de input:

\[
\texttt{WHILE X < N}
\]

este estimată ca:

\[
O(n)
\]

Prin urmare, scopul analizatorului este educațional: să ofere elevului o aproximare clară a eficienței algoritmului și să evidențieze impactul buclelor, al imbricării și al structurii codului.

\section{Concluzie}

Motorul MiniScript+ poate fi descris matematic ca un interpretor bazat pe transformări succesive ale codului: tokenizare, parsare, construire AST și evaluare. Fiecare instrucțiune modifică o stare internă, iar expresiile sunt evaluate recursiv pe baza arborelui sintactic abstract.

Această abordare permite rularea programelor, debugging pas cu pas, evidențierea liniei curente și estimarea complexității algoritmului. Astfel, MiniScript+ nu este doar un limbaj educațional simplificat, ci și un model practic pentru înțelegerea modului în care funcționează interpretoarele reale.

\end{document}

\end{document}
